\section{Data and design}\label{sec:design}

The pre-registration, its amendments and its deviations log are part of the replication package.
Appendix~\ref{app:prereg} summarizes the log (Table~\ref{tab:a-deviations}), and Appendix~\ref{app:notation}
collects the notation (Table~\ref{tab:notation}).

\subsection{Universe, sample periods and returns}\label{sec:universe}

The universe is the point-in-time membership of the S\&P 500: a stock is in the sample on a session only if it
was an index member on that session, and companies with several share classes count once. Each month the
members are split into three size tiers by market capitalization at the previous month-end: the 100 largest,
ranks 101--250 and ranks 251--500. The sample is split by time into a training period (2024-07-01 to
2025-06-30), two validation quarters and a test period. The first validation quarter, val-A (2025-07-01 to
2025-09-30), is used for search and for fitting calibration maps, temperatures and thresholds; the second, val-B
(2025-10-01 to 2025-12-31), for confirmation and the final selection. The test period, 2026-01-02 to
2026-07-31, was blinded until the analysis plan was frozen and was then used once. A target enters an
evaluation only if its whole horizon ends inside its period, and the first five sessions of each evaluation
period are left out.

Daily prices are adjusted for splits, dividends and spin-offs, and the opening price is put on the same basis.
For stock $i$ on session $t$, the close-to-close return splits at the open into an overnight return, from the
previous close to the open, and an intraday return, from the open to the close:
\begin{equation}\label{eq:split}
1 + r^{cc}_{i,t} = \bigl(1 + r^{on}_{i,t}\bigr)\bigl(1 + r^{id}_{i,t}\bigr),
\qquad
z^{k}_{i,t} = \frac{\AR^{k}_{i,t}}{\sigma^{k}_{i,t-1}}, \quad k \in \{cc, on, id\}.
\end{equation}
In logs, the close-to-close move is the sum of the overnight move and the move after the open. Each part is made
abnormal by removing the market and the stock's sector, $\AR^{k}_{i,t} = r^{k}_{i,t} - \beta^{m}_{i}\,
r^{k}_{m,t} - \beta^{s}_{i}\,(r^{k}_{s,t} - r^{k}_{m,t})$, where $r_{m}$ is the return of the S\&P 500 ETF
(SPY), $r_{s}$ the equal-weighted return of the other members of the stock's sector, and the betas come from
the 252 sessions up to $t-1$. It is then divided by $\sigma^{k}_{i,t-1}$, a robust scale (1.4826 times the
median absolute deviation) of the stock's abnormal returns over the 126 sessions up to $t-1$. A value of
$z^{k}_{i,t}$ is therefore a move in units of the stock's typical daily move, scaled by information known before
the session. Equation~\eqref{eq:split} defines what ``priced at the open'' means in this paper: news is priced
at the open when a model's reading predicts $z^{on}$ but not $z^{id}$. Where the indices are clear we write
$z_{cc}$, $z_{on}$ and $z_{id}$, as the tables and figures do. Multi-day targets add up abnormal returns over
their window and divide by the window's ex-ante volatility; the main one, $y_2$, runs from the next open to the
close of $t+5$.

\subsection{News, events and decision points}\label{sec:news}

The news comes from a commercial company-news feed (titles and snippets with publication times). An article is
available at its time stamp, or at the end of the hour or day when the stamp is rounded to the hour or carries
only a date. We drop items marked as non-news, price commentary and promotion, automatic institutional-holdings
reports, non-English articles, near-duplicates, and two kinds of item whose time stamp precedes its content:
explainers of a stock's move and running market stories updated in place. We fix three decision points per
session $t$:
\begin{itemize}[leftmargin=*, itemsep=1pt]
\item \textbf{D0}, at 09:15~ET, before the open. The model reads the news published since the previous close;
  a trade enters at the official open.
\item \textbf{D1}, at 10:00~ET, after the first 30-minute bar. The model also sees the move so far; a trade
  enters at the price at the end of that bar.
\item \textbf{D2}, after the close. The model reads the day's news and the day's move; a trade enters at the
  next open.
\end{itemize}
The observed reaction $G$ is the move up to the decision point in typical-move units: at D2 it is the day's
$z_{cc}$, and at D1 the abnormal move from the previous close to 10:00, divided by its ex-ante volatility. Its
level version $\Glvl$ divides by $\sigma^{cc}$ without a square-root-of-time adjustment, so that moves at
different horizons add up.

A learned importance filter, a logistic regression on sentence embeddings of the news \citep{xiao2024}, selects
the company-days that enter the analysis. A pre-registered rule chose it on val-A among seven candidate filters:
at a budget of 30\% of covered company-days it captures \res{f5-recall} of the company-days with a material
price reaction. An \emph{event} is a company-session with selected news: D0 events use the news up to 09:15, D2
events the news up to the close. The test period has \res{test-d0-events} D0 events on \res{test-d0-sessions}
sessions and \res{test-d2-events} D2 events.

Each event becomes a \emph{state}, a JSON document with fixed keys: the company's context (sector, size tier,
market sensitivity, volatility regime, recent abnormal performance, analyst-estimate revisions and whether an
earnings report is scheduled), the market's context (the previous session's market move, the level and change
of the VIX, rates, credit spreads, the dollar, oil and scheduled macroeconomic releases), attention (the number
of articles against the stock's usual number) and up to six news items with their source type and timing. Every
number passes through a codebook of words: a move of between one and two typical daily moves is ``moderate'',
followed by its multiple of a typical day's move. A post-reaction state
adds the reaction so far: the stock's move, the market's move over the same period and the volume against
normal.

\paragraph{The two readings.} The \emph{pre-open reading} asks, among other questions, ``Once investors fully
absorb the news, how will the company's stock perform relative to the overall market and its sector over the
next five trading days?'' on nine levels, from a very large fall (more than three typical daily moves) to a very
large rise. With $p_k$ the probability of level $k$ and $m_k$ the level's midpoint in typical daily moves, the
\emph{news-implied move} is the model's expected move,
\[
\Nhat = \sum_{k=1}^{9} p_k m_k, \qquad m = (-3.5,\ -2.5,\ -1.5,\ -0.75,\ 0,\ 0.75,\ 1.5,\ 2.5,\ 3.5).
\] The \emph{news-implied value} $V$ maps $\Nhat$ to the realized five-day level move
with an isotonic regression \citep{barlow1972} fitted on val-A for each model: it is the model's estimate of where the abnormal
price settles once the news is absorbed. The \emph{post-reaction reading} also asks whether the reaction so far
is in the wrong direction, much too small, too small, about right, too large or much too large for what the
news justifies, and whether the move will reverse, stall or continue over the next five days. Its judgment
$J = P(\text{much too small}) + P(\text{too small}) - P(\text{too large}) - P(\text{much too large})$ is positive
when the model calls the reaction an underreaction. A pre-registered check found that the post-reaction readings
anchor on the move the model is shown, so $\Nhat$ always comes from the pre-open reading.

\paragraph{Masking.} The states hide the company. Names, aliases, tickers and cashtags become ``the company'';
person names become ``an executive'' or ``a person''; dates, weekdays, years and fiscal-period labels become
``recently'' or ``the latest quarter''; share-price levels and every percent price move of any security (the
company's, other companies', indexes, commodities and currencies) are removed. Fundamental figures, such as
earnings per share, revenue and guidance, are kept, because they are the news. An independent leak scanner,
written without the masker's rules, and reading of the remaining candidates checked the result; the residues
that remain were accepted below half a percent of states for names, tickers and dates, and are disclosed in the
deviations log. Masking cannot hide a company whose news is distinctive, and Section~\ref{sec:memorization}
measures how often the models still recognize it.

\subsection{The models}\label{sec:models}

\paragraph{Laya-Fin, the primary instrument.} Laya is an open-weights decision model built on the
ModernBERT-large encoder \citep{convaiinnovations2026,warner2025}. We fine-tuned it on the training period's states with three
heads: the nine-level move (labelled with the realized five-day move), a five-level magnitude head (the size of
the same-day move) and, for the post-reaction reading, a three-class continuation head. Training ran for one
epoch on two T4 GPUs of Kaggle's free tier, with AdamW \citep{loshchilov2019} and early stopping on val-A, and one
temperature per head was fitted on val-A. The pre-registered plan called for three seeds; one seed (13) was
trained before the freeze, and the other seeds are pending robustness runs (Section~\ref{sec:limits}).

\paragraph{Commercial decision models.} Every eligible decision model in a public model router's catalogue, plus
the local Laya models, entered a two-stage zero-shot tournament on val-A: \res{tourn-models} models answered the
same 700 states. A model confidence set \citep{hansen2011} kept all of them, because about a dozen events a
session cannot separate daily ICs. The six with the highest first-stage IC then answered 3,000 states and faced
four pre-registered screens: company identification on masked states (pass at 40\% or less, against 25\% by
chance), the share of counterfactual edits (beat or miss, raise or cut, upgrade or downgrade) that move the
reading the right way (pass at 70\% or more), the IC gain from unmasking, and a shuffled-news control. No model
passed: all six identified the masked company in \res{mem-na06-min} to \res{mem-na06-max} of states, and five
also failed the counterfactual screen. By a decision logged before the freeze, the three with the highest
second-stage IC, d1-ultra, span-02 and drex-v2.0, were carried to the test as \emph{screen-failing,
model-specific comparators}: we report their results, but they cannot support a general claim. All three were
released in October 2026, after the sample ended.

\paragraph{The untuned model, the anchor and the baselines.} Laya zero-shot is the untuned Laya model answering
the same questions. The anchor is the December 2024 ModernBERT-large checkpoint \citep{warner2025}, fine-tuned
with the same heads and recipe as Laya-Fin; it was meant to show what training adds to a model whose
pretraining ended before the sample. After one epoch its outputs are nearly constant, so we report it as
\emph{uninformative}, not as a null, and make no claim about the effect of training. The baselines are vendor
sentiment and a gradient-boosted model \citep{ke2017} on the vendor's topic tags, both built on labels the vendor
produced after the fact and therefore an upper bound that may carry look-ahead, not a benchmark; a ridge
regression \citep{hoerl1970} on sentence embeddings of the news; a gradient-boosted model on price, volume and
size alone; and a generative LLM prompted for the same nine-level forecast on a 2,000-event subsample.

\subsection{Pre-registration, freeze and blinding}\label{sec:prereg}

The pre-registration was written on 2026-10-24, before any panel, target or model output existed. It fixes the
universe, the splits, the targets, the masking, the questions, the selection rules for filters, encodings and
models, the hypotheses with their families and corrections, the strategies and the inference. Three amendments
were added the same day, before any data: convergence to the news-implied value, a wider question battery and
the firm-level explanation analysis. Changes after that go only into an append-only deviations log, with the
date, the reason and whether they came before or after the freeze. Selection used validation data only.

On 2026-10-28, before any 2026 model call, the analysis plan was frozen: a SHA-256 hash over the
pre-registration (up to its deviations log), every configuration file and the code of the targets, the
reactions and the analysis specification was recorded and tagged in version control. Prices from 2026 were
sealed: they could be read only through a lock that checks the frozen hash and logs every read. The lock was
opened after the freeze to build the post-reaction states (reactions only), and once, in a single logged
sequence, to release the test outcomes. Three append-only, hash-chained ledgers record API spend (capped at
\$30; \$\res{spend-total} used), every evaluated variant as a trial (\res{ntrials-strat} when the pre-registered
strategies were evaluated, \res{ntrials-total} after the exploratory study) and every access to test-period
data. The deviations log holds \res{dev-rows} entries; Appendix~\ref{app:prereg} summarizes the material ones.

\subsection{Inference}\label{sec:inference}

The main statistic is the daily information coefficient,
\begin{equation}\label{eq:ic}
\IC_t(\Nhat, z) = \rho_{S}\bigl(\{\Nhat_{i,t}\}_{i \in \mathcal{E}_t},\, \{z_{i,t}\}_{i \in \mathcal{E}_t}\bigr),
\qquad
\overline{\IC} = \frac{1}{T}\sum_{t=1}^{T} \IC_t(\Nhat, z),
\end{equation}
the Spearman rank correlation $\rho_{S}$ between the readings and the moves across the events $\mathcal{E}_t$ of
session $t$ (sessions with at least five events), averaged over the $T$ sessions of a period. Its $t$-statistic
uses Newey--West standard errors with five lags \citep{newey1987}. A positive IC means that the readings rank
the day's moves in the right order on average; the IC is the usual measure of a cross-sectional signal. The
hypotheses that news is priced at the open compare the IC with the overnight move and with the move after the
open:
\begin{equation}\label{eq:h1}
\text{H1: } \Delta = \frac{1}{T}\sum_{t}\Bigl[\IC_t(\Nhat, z^{on}) - \IC_t(\Nhat, z^{id})\Bigr] > 0,
\qquad
\text{H1b: } \frac{1}{T}\sum_{t} \IC_t(\Nhat, z^{id}) > 0.
\end{equation}
H1 says that the news is priced mostly at the open, H1b that it is not fully priced there, so that a position
taken at the open can still earn. Both are tested one-sided with a stationary block bootstrap over sessions
\citep{politis1994}, with blocks of five sessions and 10,000 draws. H1b was added to the pre-registration before
the freeze, because the pre-registered strategy that enters at the open needs it.

The five primary hypotheses (H1, H1b, H2, H3 and H4, defined where they are tested) form one family, corrected by
Holm's method at 0.05 \citep{holm1979}; the secondary hypotheses form families corrected by
Benjamini--Hochberg at $q = 0.10$ \citep{benjamini1995}. Regressions are pooled OLS with day-by-tier fixed
effects and standard errors clustered by date and by firm \citep{cameron2011}, with Driscoll--Kraay standard
errors \citep{driscoll1998} as a check. A primary result counts as a general claim only if it holds with the same
sign and $p < 0.05$ for the primary instrument and for the anchor; if it holds for the primary instrument alone,
it is model-specific. Because the anchor is uninformative and the commercial models failed the screens, no
result in this paper can be more than model-specific.
